% !TEX root = ../main.tex
\chapter*{Abstract}

This engineering project concerns the cooperation of multiple large language models (LLMs) in solving complex problems. LLMs are usually used individually, as tools for editing texts, retrieving information and creating content. This project, however, investigated whether combining several models into a single system can yield better results than a single model working alone. The MLLMS application, developed as part of the project, is inspired by the hive mind found in nature. It was written in Python and connects through the Ollama API to locally run models. The application is operated from the terminal through a text user interface (TUI), for which no additional library was used. The models share a common space in which they exchange insights and calculation results. The responses they generate are merged into a single string, which serves as context for subsequent reasoning steps. As a result, each model has access to the conclusions of the others and can build its own answers on them. The experiments showed... On this basis, best practices for designing and implementing systems based on cooperating LLMs were formulated. The project demonstrates the potential of teamwork among language models and provides a foundation for further work on improving similar solutions.

\vspace{\baselineskip}

\noindent
\textbf{Keywords:} large language model (LLM), multi-model collaboration, hive mind, Ollama

\noindent
\textbf{Field of science and technology in accordance with OECD requirements:} Computer Science and Information technology